\documentclass[letterpaper,11pt]{moderncv}
\moderncvstyle{banking}
\moderncvcolor{black}

\usepackage[letterpaper, margin=0.8in]{geometry}
\usepackage{xcolor}
\usepackage{fontspec}
\setmainfont{Times New Roman}
\usepackage[unicode]{hyperref}

\name{Saeyoung}{Kim}
\address{Cambridge, MA, USA}{}{}
\phone[mobile]{+1 (351) 322 5510}
\email{saeyoung\_kim@uml.edu}
\social[linkedin]{saeyoung-macx-kim}

\recipient{Hiring Team}{Persimmon Technologies Corporation\\Bedford, MA}
\date{\today}
\opening{Dear Hiring Team,}
\closing{Sincerely,}

\begin{document}

\makelettertitle

I'm applying for the Manager, Robotic Manufacturing Engineering position. I've spent the last several years doing manufacturing engineering work in aerospace, and the thing I keep coming back to is wanting to own the full picture: the processes, the documentation, the tooling, the people, and how it all fits together on the production floor. That's what this role is.

At Hanwha Systems, I worked on military satellite programs where my job sat right at the boundary between engineering and production. I wrote and maintained controlled assembly and test procedures, managed the transition of design intent into production-ready documentation, and provided hands-on floor support when builds ran into problems. When hardware was discrepant, I did the root cause analysis using structured methods, drove corrective actions, and made sure the fix stuck. I also trained and mentored junior engineers and technicians on procedures and troubleshooting, which is the part of the work I found most rewarding. I contributed to facility layout planning for a new high-throughput production line, thinking about workflow, cycle time, and bottlenecks. And I managed the purchasing, integration, and commissioning of test and tooling systems.

More recently at UMass Lowell, I've been building Python-based SPC and process monitoring tools for manufacturing, including automated data pipelines, scheduled reporting, and analytics for identifying process deviations. This has given me direct experience with the data automation and process control side of manufacturing engineering.

I'll be straightforward about the gap: the role asks for 8-10 years of manufacturing experience, and my direct production floor time is closer to the Hanwha period plus my PhD hardware work. What I do bring is the full range of skills the role describes: process development, documentation systems, root cause analysis, cross-functional coordination, design review participation, tooling development, data automation, and team mentorship. I'm also someone who genuinely wants to be close to the build, not abstracted away from it.

I'm based in Cambridge, so Bedford is a short commute. My green card is currently pending, and I expect to have it in hand within the next few months. I'm happy to discuss timing directly.

\makeletterclosing

\end{document}
